\paragraph{What the two layers buy.}
The definition alone is exact and never wrong about what has happened, but before anything moves it reports every scene as one object (\accNever\% exact partitions, almost all of them things that never come apart). The memory adds hypotheses about what would come apart, learned only from what the definition recorded, and \cref{prop:identity} guarantees that a wrong hypothesis can never corrupt identity: the moment a part is moved, the definition takes over. The memory is also inspectable: every reading can be traced to the stored entries that produced it (the application shows them live), and every entry to the event that created it.

\paragraph{Carried whole is weak evidence.}
The most consequential modelling choice turned out to be how much a \emph{carried-whole} entry counts. Carrying a bookshelf carries its books, yet books come out of shelves easily: being carried along says only that a part did not come away \emph{during that motion}. The definition's own wording is careful here---a patch set is one object because nothing has been \emph{observed} to open it---and the memory must be equally careful. Chosen on validation scenes, the weight of a carried-whole entry is $w_0=\chosenW$ against $1$ for a part that came away; the ablation shows how strongly arrangements depend on it. A principled alternative is to record which part the hand acted on: a part that was pulled and did not come away is strong evidence that it is attached, while a part that was merely carried with its support is almost none. Our session sees only patches, not the hand, so it cannot make that distinction yet.

\paragraph{Limitations.}
This is a study in a controlled world, and its conclusions are bounded by it.
(1) \emph{Sensing.} Patches lie on a lattice with known correspondence, position and visibility; real perception would have to track surfaces in depth and would add noise and occlusion. Thin diagonal structures break into fragments at our patch spacing, which the definition then correctly treats as separate objects (\cref{tab:teaching}); the definition's requirement that patches be fine enough to resolve every contact is a real constraint.
(2) \emph{Interaction.} The hand is scripted; it always takes away whole units and puts them in free space. Wide drawers dropped next to each other can hide one another, which accounts for most teaching outcomes that are not as intended.
(3) \emph{Traits and candidates.} The \nTraits{} traits are designed, physically motivated descriptors, not learned features; the memory learns only how much each counts. Candidates come from straight cuts and from regions of similar signal; a part that is neither---the same colour as its surroundings and bounded by a curved seam---cannot be proposed. The same-coloured pen caps are the clearest case of this.
(4) \emph{One behaviour.} Only coming apart is recorded. The memory's format allows others (opens, bends, rolls), each with its own relevance, but we have not studied them.
(5) \emph{Scale.} Reading is a linear scan over the memory, $O(N)$ per candidate; with \nEntries{} entries a test scene (about \candPerScene{} candidates) takes \secPerScene{}\,s on one machine. Larger memories would need approximate neighbour search. With $k$ fixed the estimate is not consistent (\cref{rem:consistency}).
(6) \emph{Comparisons.} Our baselines share the candidates, traits and decision structure, which isolates the reader; they do not include pixel-based segmentation models, which take a different input and supervision. A comparison on real images and real interaction is the obvious next step and would be the real test of whether physical traits of what came apart generalise beyond a procedural world.

\paragraph{Conclusion.}
An object can be defined by observation alone, and a system holding only that definition can still learn to anticipate what will come apart: by remembering the physical traits of what it saw separate and what it saw carried whole, and reading new things off their nearest remembered neighbours under a weighting the memory works out for itself. In a controlled world spanning caps and lids, books, drawers, fruit, pencils and things that never come apart, this recovers most unseen scenes exactly while leaving identity itself untouched.
